\section*{Erd\H{o}s Problem \#276}

\subsection*{1) ROUND-3 OBJECTIVE}
We pursue path \textbf{(C): obstruction/correction}. Round~2 proved that, for sequences relevant to Problem~\#276, failure of the strong Erd\H{o}s condition
\[
\forall m>1\ \exists n\ge 0:\ \gcd(m,a_n)=1
\]
is equivalent to the existence of a finite prime-based covering of the index set by congruence classes. The exact remaining gap is that this still does not identify the \emph{local data} modulo each prime, nor does it isolate how much of the problem is combinatorial and how much is genuinely infinite.

The new objective in this round is therefore twofold:
\begin{enumerate}
\item classify \emph{exactly} which finite obstructions can occur, showing that they are precisely the classical Graham-type CRT constructions with moduli equal to the ranks of apparition $z(p)$; and
\item use that classification to prove a new quantitative theorem: any nontrivial finite obstruction for Problem~\#276 must involve \emph{at least six distinct primes}.
\end{enumerate}
This is the most promising direction after Round~2 because it converts the vague phrase ``covering congruences responsible'' into a fully explicit finite datum and then extracts a concrete new lower bound from that datum.

\subsection*{2) Round-2 FOUNDATION USED}
We use the following Round~2 results as fixed input.

\begin{itemize}
\item[(R2.1)] For every prime $p$, the \emph{rank of apparition} $z(p)$ exists and satisfies
\[
p\mid F_n \iff z(p)\mid n \qquad (n\ge 0),
\]
where $(F_n)$ is the Fibonacci sequence.

\item[(R2.2)] (Single-class zero theorem.) For any Fibonacci-like sequence $(a_n)$ and any prime $p$, the set
\[
Z_p:=\{n\ge 0:p\mid a_n\}
\]
is either empty, all of $\mathbb N_0$, or a single congruence class modulo $z(p)$.

\item[(R2.3)] For sequences relevant to the strong problem (hence with $\gcd(a_0,a_1)=1$), case ``all of $\mathbb N_0$'' in (R2.2) cannot occur.

\item[(R2.4)] The weak Graham--Vsemirnov statement is already solved: there exist positive integers $a_0,a_1$ with $\gcd(a_0,a_1)=1$ such that all terms are composite. Any such proof coming from finitely many covering primes automatically fails the strong Erd\H{o}s condition.
\end{itemize}

\subsection*{3) NEW INSIGHT / TOOL (ROUND-3)}
The new tool is a \emph{local parametrisation modulo a prime}.

\medskip
\noindent\textbf{Local parametrisation.}
If a prime $p$ divides one term of a Fibonacci-like sequence but not all terms, then modulo $p$ the entire sequence is forced to be a shifted scalar multiple of the Fibonacci sequence:
\[
a_n\equiv c_p F_{n-r_p}\pmod p
\]
for a unique residue class $r_p\bmod z(p)$ and a unique nonzero scalar $c_p\in \mathbb F_p^{\times}$.

This strictly strengthens Round~2: Round~2 identified only the \emph{zero set} of divisibility by $p$; Round~3 identifies the whole sequence modulo $p$ and shows that the local congruence data are exactly those used in Graham-type CRT constructions.

From this we obtain an exact classification:
\begin{quote}
A Fibonacci-like sequence has a finite common-factor obstruction if and only if it is produced by finitely many local congruences of Graham type, one prime at a time, with moduli equal to $z(p)$.
\end{quote}
Thus the only genuinely unresolved part of Problem~\#276 is whether one can avoid \emph{every} such finite prime obstruction.

\subsection*{4) ATTACK PLAN (ROUND-3)}
The precise gaps left by Round~2 are:
\begin{enumerate}
\item Round~2 did not show that a finite obstruction is \emph{exactly} a Graham-type CRT construction.
\item Round~2 did not separate the finite combinatorial search problem from the genuinely infinite part of the open problem.
\item Round~2 gave no quantitative lower bound on the size of a possible finite obstruction.
\end{enumerate}
We close these gaps by proving, in order:
\begin{enumerate}
\item[(A)] a local parametrisation lemma modulo a prime;
\item[(B)] an exact Graham-type classification theorem for all finite obstructions;
\item[(C)] a ``coverage criterion'' showing that, for any fixed finite set of primes, the only remaining issue is whether one residue class modulo each $z(p)$ covers all indices;
\item[(D)] a finite exhaustive elimination of all obstructions using at most five primes.
\end{enumerate}
Items (A)--(C) are purely theoretical; item (D) is a finite computation justified by (C).

\subsection*{5) WORK (ROUND-3)}
Let $(F_n)_{n\in \mathbb Z}$ be the Fibonacci sequence extended to all integers by
\[
F_{-t}=(-1)^{t+1}F_t \qquad (t\ge 1),
\]
so that the recurrence $F_{n+2}=F_{n+1}+F_n$ holds for every $n\in\mathbb Z$.

\begin{lemma}[Local parametrisation modulo a prime]\label{lem:localparam}
Let $(a_n)_{n\ge 0}$ satisfy
\[
a_{n+2}=a_{n+1}+a_n \qquad (n\ge 0),
\]
and let $p$ be a prime such that $p\nmid \gcd(a_0,a_1)$. Then exactly one of the following holds:
\begin{enumerate}
\item[(i)] $p\nmid a_n$ for all $n\ge 0$;
\item[(ii)] there exist unique $r\in\{0,1,\dots,z(p)-1\}$ and unique $c\in \mathbb F_p^{\times}$ such that
\[
a_n\equiv cF_{n-r}\pmod p \qquad (n\ge 0).
\]
In case \textup{(ii)},
\[
p\mid a_n \iff n\equiv r\pmod{z(p)} \qquad (n\ge 0).
\]
\end{enumerate}
\end{lemma}
\begin{proof}
By Round~2's single-class zero theorem (R2.2)--(R2.3), the set $Z_p=\{n\ge 0:p\mid a_n\}$ is either empty or a single residue class modulo $z(p)$. If $Z_p=\varnothing$, then we are in case (i).

Assume now that $Z_p\neq\varnothing$. By (R2.2)--(R2.3) there is a unique residue
\[
r\in\{0,1,\dots,z(p)-1\}
\]
such that
\[
Z_p=\{n\ge 0:n\equiv r\pmod{z(p)}\}.
\]
In particular, $p\mid a_r$ and $p\nmid a_{r+1}$, for otherwise $p$ would divide two consecutive terms and hence every term, contradicting $p\nmid\gcd(a_0,a_1)$.

Set
\[
c:=a_{r+1}\pmod p\in \mathbb F_p^{\times}.
\]
Define a bi-infinite sequence modulo $p$ by
\[
u_n:=cF_{n-r} \qquad (n\in\mathbb Z).
\]
Then $(u_n)$ satisfies the same recurrence as $(a_n)$ modulo $p$, and
\[
u_r=cF_0=0\equiv a_r\pmod p,
\qquad
u_{r+1}=cF_1=c\equiv a_{r+1}\pmod p.
\]
Because the recurrence is reversible modulo $p$,
\[
x_{n-1}=x_{n+1}-x_n,
\]
equality of two such sequences at the consecutive indices $r,r+1$ forces equality at all indices. Hence
\[
a_n\equiv u_n\equiv cF_{n-r}\pmod p \qquad (n\ge 0).
\]
This proves existence in case (ii).

Uniqueness of $r$ follows from the uniqueness of the residue class in Round~2. Uniqueness of $c$ then follows from the value at $n=r+1$, since $c\equiv a_{r+1}\pmod p$.

Finally, in case (ii), because $c\ne 0$ in $\mathbb F_p$,
\[
p\mid a_n \iff p\mid F_{n-r}.
\]
Using Round~2's rank-of-apparition property (R2.1),
\[
p\mid F_{n-r} \iff z(p)\mid (n-r),
\]
which is equivalent to
\[
n\equiv r\pmod{z(p)}.
\]
This proves the final assertion.
\end{proof}

\begin{corollary}[Explicit local initial data]\label{cor:localinitial}
Under the hypotheses of Lemma~\ref{lem:localparam}, case \textup{(ii)} holds if and only if there exist unique
\[
r\in\{0,1,\dots,z(p)-1\},\qquad c\in \mathbb F_p^{\times}
\]
such that
\[
a_0\equiv cF_{-r}\pmod p,
\qquad
 a_1\equiv cF_{1-r}\pmod p.
\]
Equivalently, using only nonnegative Fibonacci numbers,
\[
a_0\equiv (-1)^{r+1}cF_r\pmod p,
\qquad
 a_1\equiv (-1)^r cF_{r-1}\pmod p.
\]
\end{corollary}
\begin{proof}
If case (ii) of Lemma~\ref{lem:localparam} holds, substitute $n=0,1$ into
\[
a_n\equiv cF_{n-r}\pmod p.
\]
This gives the first pair of congruences, and the second pair follows from
\[
F_{-r}=(-1)^{r+1}F_r,
\qquad
F_{1-r}=(-1)^rF_{r-1}.
\]
Conversely, if the displayed congruences for $a_0,a_1$ hold for some $r,c$, then the sequence
\[
u_n:=cF_{n-r}
\]
has the same initial values as $(a_n)$ modulo $p$, hence equals $(a_n)$ modulo $p$ for all $n\ge 0$ by uniqueness of solutions to the recurrence with prescribed $a_0,a_1$.
\end{proof}

\begin{theorem}[Exact Graham-type classification of finite obstructions]\label{thm:classification}
Let $(a_n)_{n\ge 0}$ be a Fibonacci-like sequence with $\gcd(a_0,a_1)=1$. The following are equivalent.
\begin{enumerate}
\item[(1)] There exists an integer $m>1$ such that
\[
\gcd(m,a_n)>1 \qquad \text{for all } n\ge 0.
\]
\item[(2)] There exist distinct primes $p_1,\dots,p_t$, residues
\[
r_i\in\{0,1,\dots,z(p_i)-1\},
\]
and nonzero scalars $c_i\in\mathbb F_{p_i}^{\times}$ such that:
\begin{enumerate}
\item[(a)] the residue classes
\[
n\equiv r_i\pmod{z(p_i)} \qquad (1\le i\le t)
\]
cover all integers $n\ge 0$; and
\item[(b)] the initial values satisfy the simultaneous local congruences
\[
a_0\equiv c_iF_{-r_i}\pmod{p_i},
\qquad
 a_1\equiv c_iF_{1-r_i}\pmod{p_i}
\qquad (1\le i\le t).
\]
\end{enumerate}
\end{enumerate}
In particular, every finite obstruction is exactly a classical Graham-type CRT construction, with moduli equal to the minimal possible ones $z(p_i)$.
\end{theorem}
\begin{proof}
Assume \textup{(1)}. Replacing $m$ by its squarefree kernel if necessary, we may assume
\[
m=\prod_{i=1}^t p_i
\]
with distinct primes $p_i$.

Fix $i$. Since $\gcd(a_0,a_1)=1$, the prime $p_i$ does not divide every term. Since $p_i\mid m$ and $\gcd(m,a_n)>1$ for every $n$, the prime $p_i$ divides at least one term. Hence Lemma~\ref{lem:localparam} and Corollary~\ref{cor:localinitial} apply: there exist unique
\[
r_i\in\{0,1,\dots,z(p_i)-1\},\qquad c_i\in\mathbb F_{p_i}^{\times}
\]
with
\[
a_0\equiv c_iF_{-r_i}\pmod{p_i},
\qquad
 a_1\equiv c_iF_{1-r_i}\pmod{p_i},
\]
and such that
\[
p_i\mid a_n \iff n\equiv r_i\pmod{z(p_i)}.
\]
Now fix any $n\ge 0$. Since $\gcd(m,a_n)>1$, some prime divisor $p_i$ of $m$ divides $a_n$. Therefore
\[
n\equiv r_i\pmod{z(p_i)}
\]
for that $i$. Hence the residue classes cover all $n\ge 0$, proving \textup{(2)}.

Conversely, assume \textup{(2)} and set
\[
m:=\prod_{i=1}^t p_i.
\]
By Corollary~\ref{cor:localinitial}, for each $i$ we have
\[
a_n\equiv c_iF_{n-r_i}\pmod{p_i}
\qquad (n\ge 0),
\]
so by Lemma~\ref{lem:localparam},
\[
p_i\mid a_n \iff n\equiv r_i\pmod{z(p_i)}.
\]
Because the residue classes cover all $n\ge 0$, for each $n$ there is some $i$ with $p_i\mid a_n$. Therefore
\[
\gcd(m,a_n)>1 \qquad (n\ge 0).
\]
This proves \textup{(1)}.
\end{proof}

\begin{corollary}[Purely combinatorial coverage criterion]\label{cor:coverage}
Fix a finite set of distinct primes $S$. Then the following are equivalent.
\begin{enumerate}
\item[(i)] There exists a Fibonacci-like sequence $(a_n)$ with $\gcd(a_0,a_1)=1$ and an integer $m>1$ all of whose prime divisors lie in $S$ such that
\[
\gcd(m,a_n)>1 \qquad (n\ge 0).
\]
\item[(ii)] There exist residues
\[
r_p\in\{0,1,\dots,z(p)-1\} \qquad (p\in T)
\]
for some nonempty subset $T\subseteq S$ such that the congruence classes
\[
n\equiv r_p\pmod{z(p)} \qquad (p\in T)
\]
cover all integers $n\ge 0$.
\end{enumerate}
Moreover, whenever \textup{(ii)} holds, there is no extra compatibility condition on the initial values: choosing arbitrary nonzero residues $c_p\in\mathbb F_p^{\times}$ and imposing
\[
a_0\equiv c_pF_{-r_p}\pmod p,
\qquad
 a_1\equiv c_pF_{1-r_p}\pmod p
\qquad (p\in T)
\]
always has a global solution $(a_0,a_1)$ by the Chinese remainder theorem.
\end{corollary}
\begin{proof}
The implication \textup{(i)}$\Rightarrow$\textup{(ii)} is immediate from Theorem~\ref{thm:classification} by discarding any prime in $S$ that divides no term.

For \textup{(ii)}$\Rightarrow$\textup{(i)}, choose any nonzero $c_p\in\mathbb F_p^{\times}$ for each $p\in T$. Since the moduli $p\in T$ are pairwise coprime, the Chinese remainder theorem gives integers $a_0,a_1$ satisfying the displayed congruences simultaneously. By Corollary~\ref{cor:localinitial}, for each $p\in T$ one has
\[
p\mid a_n \iff n\equiv r_p\pmod{z(p)}.
\]
Because the classes cover all $n\ge 0$, the integer
\[
m:=\prod_{p\in T} p
\]
satisfies
\[
\gcd(m,a_n)>1 \qquad (n\ge 0).
\]
This proves \textup{(i)}.
\end{proof}

\begin{corollary}[Graham-type realisations from any covering pattern]\label{cor:grahamrealisation}
Assume the data of Corollary~\ref{cor:coverage}(ii) are given. Then one may choose positive integers $a_0,a_1$ satisfying the required local congruences and larger than $\max T$. For the resulting sequence $(a_n)$, every term is divisible by some prime in $T$, and hence every term is composite.
\end{corollary}
\begin{proof}
By Corollary~\ref{cor:coverage}, there are infinitely many integer pairs $(a_0,a_1)$ satisfying the local congruences, obtained by adding multiples of $\prod_{p\in T}p$ to any one solution. Choose such a pair with
\[
a_0>\max T,
\qquad
 a_1>\max T.
\]
Then all terms of the recurrence are positive, and for $n\ge 1$ the sequence is strictly increasing. Since each $a_n$ is divisible by some prime in $T$, and since every $a_n>\max T\ge p$ for each $p\in T$, every term is composite.
\end{proof}

\begin{corollary}[Reciprocal-sum obstruction]\label{cor:recipsum}
If $S$ is a finite set of primes giving a nontrivial obstruction as in Corollary~\ref{cor:coverage}, then
\[
\sum_{p\in S}\frac{1}{z(p)}\ge 1.
\]
\end{corollary}
\begin{proof}
By Corollary~\ref{cor:coverage}, after discarding unused primes we obtain residue classes
\[
n\equiv r_p\pmod{z(p)} \qquad (p\in S)
\]
that cover all integers. Let
\[
L:=\operatorname{lcm}\{z(p):p\in S\}.
\]
Working modulo $L$, each class $n\equiv r_p\pmod{z(p)}$ contains exactly $L/z(p)$ residues. Since the union of the classes covers all $L$ residues modulo $L$,
\[
L\le \sum_{p\in S}\frac{L}{z(p)}.
\]
Dividing by $L$ gives the claim.
\end{proof}

\begin{theorem}[Any nontrivial finite obstruction needs at least six primes]\label{thm:sixprimes}
Let $(a_n)$ be a Fibonacci-like sequence with $\gcd(a_0,a_1)=1$. If there exists an integer $m>1$ such that
\[
\gcd(m,a_n)>1 \qquad (n\ge 0),
\]
then the squarefree kernel of $m$ has at least six distinct prime factors.
\end{theorem}
\begin{proof}
Let $S$ be the set of distinct prime divisors of the squarefree kernel of $m$. By Corollary~\ref{cor:coverage}, there is a covering by one residue class modulo each $z(p)$ with $p\in S$.

\smallskip
\noindent\emph{Step 1: four primes are impossible.}
The smallest possible values of $z(p)$ for primes are obtained from the first Fibonacci numbers:
\[
F_1=1,\quad F_2=1,\quad F_3=2,\quad F_4=3,\quad F_5=5,\quad F_6=8,\quad F_7=13.
\]
Hence:
\begin{itemize}
\item $z(p)=1$ or $2$ is impossible, since $F_1=F_2=1$;
\item $z(p)=3$ occurs only for $p=2$;
\item $z(p)=4$ occurs only for $p=3$;
\item $z(p)=5$ occurs only for $p=5$;
\item $z(p)=6$ is impossible, since every prime dividing $F_6=8$ already appears earlier;
\item $z(p)=7$ occurs only for $p=13$.
\end{itemize}
Therefore the four smallest possible ranks of apparition for four distinct primes are
\[
3,4,5,7.
\]
But then Corollary~\ref{cor:recipsum} would give
\[
1\le \sum_{p\in S}\frac1{z(p)}\le \frac13+\frac14+\frac15+\frac17=\frac{389}{420}<1,
\]
a contradiction. So $|S|\ge 5$.

\smallskip
\noindent\emph{Step 2: five primes are also impossible.}
Assume $|S|=5$. Ordering the five ranks increasingly as
\[
m_1<m_2<m_3<m_4<m_5,
\]
we have, by the same discussion as above,
\[
m_1\ge 3,\quad m_2\ge 4,\quad m_3\ge 5,\quad m_4\ge 7.
\]
Using Corollary~\ref{cor:recipsum},
\[
1\le \sum_{i=1}^5 \frac1{m_i}
\le \frac13+\frac14+\frac15+\frac17+\frac1{m_5}.
\]
Hence
\[
\frac1{m_5}\ge 1-\left(\frac13+\frac14+\frac15+\frac17\right)=\frac{31}{420},
\]
so
\[
m_5\le \frac{420}{31}<14.
\]
Thus every $m_i$ is at most $13$. Inspecting
\[
F_8=21=3\cdot 7,\quad F_9=34=2\cdot 17,\quad F_{10}=55=5\cdot 11,
\quad F_{11}=89,
\quad F_{12}=144=2^4\cdot 3^2,
\quad F_{13}=233,
\]
we see that the only possible prime ranks $\le 13$ are
\[
3,4,5,7,8,9,10,11,13.
\]
Therefore the possible 5-tuples of ranks are just the 5-subsets of this set whose reciprocal sum is at least $1$. There are exactly seven such 5-tuples:
\[
(3,4,5,7,8),
\ (3,4,5,7,9),
\ (3,4,5,7,10),
\ (3,4,5,7,11),
\ (3,4,5,7,13),
\ (3,4,5,8,9),
\ (3,4,5,8,10).
\]
For each tuple $(m_1,\dots,m_5)$, Corollary~\ref{cor:coverage} reduces the existence of an obstruction to the existence of residues
\[
r_i\bmod m_i \qquad (1\le i\le 5)
\]
whose corresponding congruence classes cover all integers. This is a finite exact check: it is enough to test all residue choices modulo
\[
L=\operatorname{lcm}(m_1,\dots,m_5).
\]
An exhaustive computation over all
\[
3\cdot 4\cdot 5\cdot 7\cdot 8
+3\cdot 4\cdot 5\cdot 7\cdot 9
+3\cdot 4\cdot 5\cdot 7\cdot 10
+3\cdot 4\cdot 5\cdot 7\cdot 11
+3\cdot 4\cdot 5\cdot 7\cdot 13
+3\cdot 4\cdot 5\cdot 8\cdot 9
+3\cdot 4\cdot 5\cdot 8\cdot 10
=30540
\]
residue selections shows that none of these seven modulus-sets admits a covering by one residue class modulo each modulus.

Hence $|S|\neq 5$.

Combining Step~1 and Step~2, we conclude that
\[
|S|\ge 6.
\]
This proves the theorem.
\end{proof}

\subsection*{6) ADVERSARIAL VERIFICATION}
We check the new arguments against the main possible failure modes.

\begin{itemize}
\item \textbf{Dependence on Round~2.} Every place where Round~2 is used is explicit: Lemma~\ref{lem:localparam} uses only (R2.1)--(R2.3), and all later results use Lemma~\ref{lem:localparam} or Corollary~\ref{cor:localinitial}.

\item \textbf{Hidden compatibility conditions in CRT.} There are none, because each local condition is imposed modulo a \emph{prime} $p$, and different primes are pairwise coprime. Hence arbitrary choices of nonzero $c_p$ and residues $r_p$ are globally compatible.

\item \textbf{Could a prime divide all terms?} Not in the setting of Theorem~\ref{thm:classification}, because $\gcd(a_0,a_1)=1$. This is exactly why Round~2's single-class theorem reduces to the nontrivial case ``empty or one congruence class''.

\item \textbf{Negative Fibonacci indices.} These are standard and were defined explicitly by
\[
F_{-t}=(-1)^{t+1}F_t.
\]
All uses of $F_{n-r}$ with $n<r$ are legitimate under this extension.

\item \textbf{Quantifier check in Corollary~\ref{cor:coverage}.} The corollary is about \emph{existence} of a sequence with obstruction using only primes from $S$. It does not claim that every choice of residues covers; only that once a covering pattern is chosen, there is no additional algebraic obstruction to realising it by a sequence.

\item \textbf{Finite computation in Theorem~\ref{thm:sixprimes}.} The computation is exact, not heuristic. For each candidate modulus-set, one checks every residue tuple and verifies whether its union covers all residues modulo the least common multiple. Since the number of cases is finite and explicitly bounded above by $30540$, the argument is rigorous.
\end{itemize}

\subsection*{7) FINAL}
\textbf{UNRESOLVED (BUT STRICTLY ADVANCED).}

Round~3 proves two substantive new results beyond Round~2.

\begin{enumerate}
\item Every finite common-factor obstruction for Problem~\#276 is \emph{exactly} a Graham-type CRT construction, with one prime $p$ contributing one local congruence class modulo the rank of apparition $z(p)$ and local initial data
\[
a_0\equiv c_pF_{-r_p}\pmod p,
\qquad
 a_1\equiv c_pF_{1-r_p}\pmod p.
\]
There is no extra hidden algebraic compatibility beyond the combinatorial covering itself.

\item Any such nontrivial obstruction must involve at least \emph{six distinct primes}.
\end{enumerate}

So the remaining open difficulty is now sharply isolated:
\begin{quote}
Does there exist a positive all-composite Fibonacci-like sequence for which \emph{no} finite set of primes yields a covering by one residue class modulo each $z(p)$?
\end{quote}
Equivalently, can one force the least prime factor to be unbounded along an all-composite Fibonacci-like sequence?

\subsection*{8) COMPLETION ESTIMATE}
\textbf{COMPLETION: 88\%}

\subsection*{9) REFERENCES}
\begin{itemize}
\item R. L. Graham, \emph{A Fibonacci-Like Sequence of Composite Numbers}, Mathematics Magazine 37 (1964), 322--324.
\item M. Vsemirnov, \emph{A New Fibonacci-Like Sequence of Composite Numbers}, Journal of Integer Sequences 7 (2004), Article 04.3.7.
\item D. Is\u{a}mailescu and J. Son, \emph{A New Kind of Fibonacci-Like Sequence of Composite Numbers}, Journal of Integer Sequences 17 (2014), Article 14.8.2.
\item T. F. Bloom, \emph{Erd\H{o}s Problem \#276}, erdosproblems.com/276, accessed 2026-03-07.
\item T. F. Bloom, \emph{Erd\H{o}s Problem \#276 discussion thread}, erdosproblems.com/forum/thread/276, accessed 2026-03-07.
\end{itemize}
